\chapter{Auctions and the Close}\label{ch:hf:auctions-and-the-close}

At ten to four the exchange publishes a buy imbalance of two hundred thousand shares in a stock. In the next ten minutes someone will sell them
into the auction, and the price at which they do is set by whoever shows the best offer. In this chapter's closing book the imbalance alone would
lift the close by 14 cents on a \$100 stock; a provider that offers 120,000 shares at the reference price holds the move to 9 cents and earns 6.9
cents a share, if 85\% of the move would have reverted by the next morning. Late orders that change the imbalance cut that to 5.5 cents, unless the
provider asks for a premium.

\section{The closing auction as a liquidity event}

The closing auction (Book 1, chapter 13) sets the price at which index funds, benchmarked managers and derivatives settle. That makes it the day's
largest single pool of liquidity, and one whose buyers and sellers care about getting the closing price more than about the price itself.
Bogousslavsky and Muravyev measured the closing auction at 7.5\% of daily US volume in 2018, up from 3.1\% in 2010. Closing prices typically matched
the pre-close bid or ask; price impact was lower than in continuous trading; and auction price deviations reverted quickly and almost completely, on
average. Auction volume relative to intraday volume jumped on S\&P 500 additions and rose permanently afterwards, which they read as the growth of
indexing and exchange-traded funds feeding the close.

A market maker's role there is to provide the other side of an imbalance, at a price, and carry the position overnight or hedge it.

\section{Reading the imbalance publication}

Before the auction the exchange publishes what it would do if it uncrossed now: the reference or indicative price, the paired shares, the
imbalance and its side (the imbalance publication of One Quant Book 10, chapter 10).

\begin{dated}{2026-09}{One exchange's closing schedule}\label{dat:hf:auctions-and-the-close:nasdaq}
Nasdaq's frequently asked questions on its crosses (as published in November 2025) state that between 3:50 and 3:55 p.m. it disseminates the current
reference price, the paired shares, the imbalance shares and the imbalance side every 10 seconds, and between 3:55 and 4:00 p.m. every second, adding
near and far indicative clearing prices. Market-on-close orders must be received before 3:55 p.m.; limit-on-close orders before 3:58 p.m., those
after 3:55 re-priced if more aggressive than the 3:50 or 3:55 reference prices; on-close orders can be modified or cancelled before 3:50 p.m.
Imbalance-only orders exist to offset a market-on-close imbalance.
\end{dated}

The publications are a forecast, not the answer: market-on-close orders keep arriving until their cutoff and limit-on-close orders after it. A
provider that sizes its order on a publication must allow for the imbalance changing before the uncross.

\section{Offsetting imbalances: pricing the order}

\begin{definition}[Imbalance-offsetting order]\label{def:hf:auctions-and-the-close:offset}
An \emph{imbalance-offsetting order}\index{imbalance-offsetting order} is a closing-auction order on the opposite side of a published imbalance,
priced and sized to trade only if the auction needs it: a sell against a buy imbalance, a buy against a sell imbalance.
\end{definition}

\begin{definition}[Auction impact curve]\label{def:hf:auctions-and-the-close:impact}
An \emph{auction impact curve}\index{auction impact curve} is the uncrossing price's move from the reference price as a function of the size of an
imbalance, given the rest of the auction book; it is the auction's supply curve seen from the imbalance.
\end{definition}

In the chapter's closing book (300,000 shares paired, limit-on-close orders of $2{,}000\,j$ shares $j$ cents from \$100 on each side), the impact
curve is concave: 5 cents for 25,000 shares, 10 for 100,000, 14 for 200,000, 20 for 400,000. The provider's offer moves the close back toward the
reference; its profit is the gap between the close and tomorrow's value on each share it sells. If a share $\pi$ of the closing move is information
that stays, and the rest reverts, the provider sells at the close $p$ and expects tomorrow's value $P_0+\pi\,\Delta(I)$, where $\Delta(I)$ is the move
the imbalance alone would have caused (\cref{lst:hf:auctions-and-the-close:provider}). The best offer trades off size against price: selling more
lowers the close for every share sold.

\begin{omfigure}
\begin{tikzpicture}
\begin{axis}[width=11cm,height=5.4cm,xlabel={\small imbalance (thousand shares)},ylabel={\small cents},xmin=0,xmax=420,ymin=0,ymax=22,
  tick label style={font=\footnotesize},grid=major,grid style={gray!25},table/col sep=comma,
  legend style={font=\scriptsize,at={(0.5,-0.3)},anchor=north},legend columns=3]
\addplot[black,thick,mark=*,mark size=1.2pt] table[x=imbalance,y=move]{figdata/hft/16-auctions-and-the-close/provider.csv};
\addplot[black,thin,dashed,mark=o,mark size=1.2pt] table[x=imbalance,y=close15]{figdata/hft/16-auctions-and-the-close/provider.csv};
\addplot[omThm,thick,mark=square*,mark size=1.2pt] table[x=imbalance,y=ps0]{figdata/hft/16-auctions-and-the-close/provider.csv};
\addplot[omDef,thick,mark=triangle*,mark size=1.4pt] table[x=imbalance,y=ps15]{figdata/hft/16-auctions-and-the-close/provider.csv};
\addplot[gray,thick,mark=diamond*,mark size=1.4pt] table[x=imbalance,y=ps50]{figdata/hft/16-auctions-and-the-close/provider.csv};
\legend{{closing move, alone},{closing move, with provider},{profit a share, all reverts},{profit a share, 15\% stays},{profit a share, 50\% stays}}
\end{axis}
\end{tikzpicture}

\omcaption{The \omterm{def:hf:auctions-and-the-close:impact}{auction impact curve} of the chapter's closing book (the close's move above \$100 against a buy imbalance), the move when the provider
makes its best offer (15\% of the move stays overnight), and the provider's expected profit a share at its best offer when none, 15\% or half of the
move stays. Uncrossed with \texttt{firm.auction}. Data: \texttt{hf\_auction.offers}.}\label{fig:hf:auctions-and-the-close:provider}
\end{omfigure}

The profit a share grows with the imbalance, less than proportionally (\cref{fig:hf:auctions-and-the-close:provider}): 4.5 cents at 100,000 shares,
6.9 at 200,000 and 10 at 400,000 when 15\% of the move stays. It depends as much on the reversal assumption as on the imbalance: at 200,000 shares, 8
cents if everything reverts, 4 if half stays. The reversal is what the provider must estimate from history, stock by stock, and it is the part of
its model that can be wrong on the day that matters: an imbalance caused by news.

\section{The auction and the continuous book before it}

The auction does not happen in isolation. Traders who see the imbalance publication can trade in the continuous book before the close, and the
closing book's limit orders can be priced from it. A provider can hedge before the close (sell some of the stock in the continuous book, or a
correlated future) or after (the closing price is a known number; the futures market trades on). The trade-at-settlement futures of Book 1,
chapter 21 let a provider lock in a hedge at the index's settlement price, removing the market part of its overnight risk at a known basis.

Late orders are the provider's main risk inside the auction. With a late-order uncertainty of 60,000 shares on a published 200,000, the offer
that was best on the publication (120,000 shares at the reference price) earns 5.5 cents a share instead of 6.9; the best offer given the
uncertainty keeps the size and asks 4 cents above the reference, earning 6.0: a premium that protects the provider when the imbalance shrinks and
its offer would otherwise set the close.

\section{Opening auctions and reopenings}

The opening auction and the reopening after a halt have the same mechanics with less information and more risk: the overnight news is not yet
in any price, the imbalance is smaller relative to the day's uncertainty, and the reference price is stale. Providers ask for more, and hedge in the index future as soon as the stock opens.

\section{Strategy files}

\begin{strategyfile}{Closing-auction imbalance liquidity provision}\label{strat:hf:auctions-and-the-close:close}
\sfield{Who pays you, and why} Index funds and benchmarked managers who must trade at the close and pay the auction's impact to do so.
\sfield{Instruments and venues} Closing auctions of the listing exchanges; imbalance-only or limit-on-close orders.
\sfield{Signal} The published imbalance and reference price, the stock's closing book, the estimated share of the move that stays overnight.
\sfield{Sizing and execution} Offer the size and price that maximise expected profit against the impact curve; add a premium for late orders; hedge
the market part with index futures at settlement.
\sfield{Costs} Overnight risk of the part that stays; exchange fees; hedge costs.
\sfield{How it dies} Information: an imbalance caused by news leaves most of the move in place (4 cents a share at half staying instead of 6.9);
competition among providers.
\sfield{Horizon, capacity, infrastructure} Ten minutes into the close and overnight; capacity from imbalances across thousands of stocks; the
imbalance feeds of every exchange.
\sfield{Backtest honestly} Publications as disseminated at each time, late orders, the uncross rules, and the next open as the exit.
\sfield{Sources} Bogousslavsky and Muravyev (2023); the dated box.
\end{strategyfile}

\begin{strategyfile}{Opening-auction liquidity provision}\label{strat:hf:auctions-and-the-close:open}
\sfield{Who pays you, and why} Traders who must execute at the open after overnight news or flows.
\sfield{Instruments and venues} Opening auctions; imbalance publications before the open.
\sfield{Signal} The opening imbalance and indicative price against a fair value from overnight futures and news.
\sfield{Sizing and execution} Smaller sizes and wider premiums than at the close; hedge in the future as soon as the stock opens.
\sfield{Costs} Overnight information not in the reference price; the first minutes' volatility.
\sfield{How it dies} News that the provider's fair value misses.
\sfield{Horizon, capacity, infrastructure} Minutes; the pre-open data of every venue.
\sfield{Backtest honestly} Futures and news as known at each pre-open time.
\sfield{Sources} Book 1, chapter 13; this chapter's closing-auction model, applied with less information.
\end{strategyfile}

\begin{strategyfile}{Close hedged with trade-at-settlement futures}\label{strat:hf:auctions-and-the-close:tas}
\sfield{Who pays you, and why} As the closing provision, with the market part of the overnight risk transferred at a known basis.
\sfield{Instruments and venues} Closing auctions of index constituents; trade-at-settlement index futures (Book 1, chapter 21).
\sfield{Signal} Aggregate imbalances across the index's constituents, which predict the index's closing move.
\sfield{Sizing and execution} Provide against constituents' imbalances; buy or sell trade-at-settlement futures for the aggregate beta during the
day, before the close is known.
\sfield{Costs} The trade-at-settlement basis, fees, the idiosyncratic part left unhedged.
\sfield{How it dies} Crowding in the trade-at-settlement basis; imbalances offsetting across names so the hedge overshoots.
\sfield{Horizon, capacity, infrastructure} The last hour and overnight.
\sfield{Backtest honestly} The trade-at-settlement prices available at each time, not the settlement itself.
\sfield{Sources} Book 1, chapter 21.
\end{strategyfile}

\section{Tutorial: two hundred thousand to sell}\label{tut:hf:auctions-and-the-close:tut}

\begin{tutorial}
\textbf{Goal.} Price an \omterm{def:hf:auctions-and-the-close:offset}{imbalance-offsetting order} from the \omterm{def:hf:auctions-and-the-close:impact}{auction impact curve} and the overnight reversal, and measure the cost of late orders.
\textbf{End state:} \cref{fig:hf:auctions-and-the-close:provider} and the numbers in the text.
\begin{tutsteps}
\item \textbf{The closing book} and its orders for \texttt{firm.auction.uncross}.
\omcode{code/firm/auctionmm/firm_auctionmm.py}{35}{50}{Paired market-on-close interest, limit-on-close orders deepening away from the reference, the imbalance and the provider.}\label{lst:hf:auctions-and-the-close:book}
\item \textbf{The provider's P\&L} against tomorrow's value, with late orders.
\omcode{code/firm/auctionmm/firm_auctionmm.py}{63}{80}{Uncross with and without the provider; tomorrow's value keeps a share of the move the imbalance alone would cause.}\label{lst:hf:auctions-and-the-close:provider}
\item \textbf{The best offer} over sizes of 20,000 to 400,000 shares and limits 0 to 10 cents through the reference (\texttt{hf\_auction.offers}).
\item \textbf{Late orders}: a standard deviation of 60,000 shares on a published 200,000 (\texttt{hf\_auction.late}).
\end{tutsteps}
\textbf{What to change next.} Add a second provider and find the equilibrium; estimate the share that stays from the history of each stock's
closing deviations; run the auction on \texttt{firm.exchsim}'s \texttt{auction\_venue} preset with its publications.
\end{tutorial}

\section{Build: the auction liquidity module}\label{bld:hf:auctions-and-the-close:auctionmm}

\begin{build}
\bfield{Purpose} Model closing books, impact curves and \omterm{def:hf:auctions-and-the-close:offset}{imbalance-offsetting orders}, with late orders and the overnight reversal.
\bfield{Interface} \texttt{ClosingBook(p0, levels, d0, paired)} with \texttt{orders(imbalance, q, k)}, \texttt{close\_price}, \texttt{impact\_curve},
\texttt{provider(book, imbalance, q, k, perm, late\_sd, n, seed)}, \texttt{best\_offer}, \texttt{publications}. Built on \texttt{firm.auction}
(Book 1).
\bfield{Rules} Prices in ticks; the provider's order is a limit $k$ ticks through the reference on the side opposite the imbalance; P\&L against
tomorrow's value.
\bfield{Acceptance tests} \texttt{code/firm/auctionmm/tests/}: a balanced book closes at the reference; an imbalance of 30 in a toy book closes two
ticks away on either side; the impact curve rises; the provider moves the close and fills; its P\&L falls when the move stays; the best offer is at
least as good as a given one; publications converge to the final imbalance.
\bfield{Stretch} Several providers; hedging with trade-at-settlement futures; opening auctions with overnight information.
\end{build}

\begin{omsources}
\item V.~Bogousslavsky, D.~Muravyev, Who trades at the close? Implications for price discovery and liquidity, \emph{Journal of Financial
Markets} 66, 2023, 100852.
\item Nasdaq, The Nasdaq Opening and Closing Crosses: frequently asked questions (November 2025).
\end{omsources}

\section{Exercises}

\begin{exercise}[$\star$]\label{exo:hf:auctions-and-the-close:1}
In a toy closing book, sell limit-on-close orders rest at 101 (10 shares), 102 (20) and 103 (30) against a reference of 100. Where does a market-on-close
buy imbalance of 30 shares close the auction?
\end{exercise}

\begin{exercise}[$\star$]\label{exo:hf:auctions-and-the-close:2}
The close moves 9 cents with the provider's order; tomorrow's value keeps 15\% of the 14-cent move the imbalance alone would cause. What does the
provider earn a share, and on 120,000 shares?
\end{exercise}

\begin{exercise}[$\star$]\label{exo:hf:auctions-and-the-close:3}
Under the dated box's schedule, until when can a market-on-close order be entered, and when does the imbalance start being published every second?
\end{exercise}

\begin{exercise}[$\star\star$]\label{exo:hf:auctions-and-the-close:4}
Why does the provider's best offer rarely fill the whole imbalance?
\end{exercise}

\begin{exercise}[$\star\star$]\label{exo:hf:auctions-and-the-close:5}
Why does the profit a share fall from 8 to 4 cents at 200,000 shares when half the move stays instead of none?
\end{exercise}

\begin{exercise}[$\star\star$]\label{exo:hf:auctions-and-the-close:6}
Why does a premium of 4 cents help when late orders make the imbalance uncertain?
\end{exercise}

\begin{exercise}[$\star\star\star$]\label{exo:hf:auctions-and-the-close:7}
\emph{Coding.} Double the depth of the closing book (\texttt{d0=4000}). What happens to the impact of a 200,000-share imbalance and to the
provider's best profit a share with 15\% staying?
\end{exercise}

\begin{exercise}[$\star\star\star$]\label{exo:hf:auctions-and-the-close:8}
\emph{Find the flaw.} ``Closing deviations revert almost completely on average, so we offset every imbalance at full size.''
\end{exercise}

\section{Problem: Two Hundred Thousand to Sell}

\begin{problem}[{Weekend problem --- two hundred thousand to sell}]\label{pb:hf:auctions-and-the-close:1}
A provider decides at 3:55 p.m. how many shares to offer against a published buy imbalance of 200,000.

\medskip\noindent\textbf{Part I --- The auction.}
\begin{enumerate}
\item Why is the close the day's largest pool of liquidity?
\item What did Bogousslavsky and Muravyev find about volume, prices and reversals at the close?
\item Summarise the dated box's schedule.
\item What does an imbalance publication contain?
\end{enumerate}

\medskip\noindent\textbf{Part II --- Pricing.}
\begin{enumerate}[resume]
\item Define an \omterm{def:hf:auctions-and-the-close:offset}{imbalance-offsetting order}.
\item Define the \omterm{def:hf:auctions-and-the-close:impact}{auction impact curve} and give its values in the chapter's book.
\item Write the provider's profit given the share $\pi$ that stays.
\item Why does selling more lower the profit a share?
\end{enumerate}

\medskip\noindent\textbf{Part III --- Measurements.}
\begin{enumerate}[resume]
\item Give the best offer and profit at 200,000 shares with 15\% staying.
\item How does the profit a share vary with the imbalance and with $\pi$?
\item What do late orders cost, and what recovers part of it?
\item How much does the provider move the close?
\end{enumerate}

\medskip\noindent\textbf{Part IV --- The verdict.}
\begin{enumerate}[resume]
\item State the \emph{named result}: the offsetting provider's expected profit a share against the imbalance size, and the share of the closing
move that reverts overnight.
\item How would you estimate $\pi$ stock by stock?
\item How would you hedge the overnight position?
\item What changes at the open?
\item What does a second provider do to the first's profit?
\item What does the provider contribute to the closing price's quality?
\item Which strategy file is most exposed to news?
\item In one sentence: what does an imbalance provider sell?
\end{enumerate}
\end{problem}

\section{Interview questions}

\begin{interviewq}[$\star$ \iqroles{trader}]\label{iq:hf:auctions-and-the-close:1}
A stock shows a 500,000-share buy imbalance at 3:50 p.m. What do you look at before offering?
\end{interviewq}

\begin{interviewq}[$\star\star$ \iqroles{researcher}]\label{iq:hf:auctions-and-the-close:2}
How would you measure the share of closing-auction price moves that reverts overnight?
\end{interviewq}

\begin{interviewq}[$\star\star$ \iqroles{developer}]\label{iq:hf:auctions-and-the-close:3}
Implement the uncrossing rule of a call auction. How do you handle ties?
\end{interviewq}

\begin{interviewq}[$\star\star$ \iqroles{risk}]\label{iq:hf:auctions-and-the-close:4}
Your desk offsets imbalances in 800 stocks each day. What is the firm's overnight risk, and how is it limited?
\end{interviewq}

\begin{interviewq}[$\star\star$ \iqroles{trader}]\label{iq:hf:auctions-and-the-close:5}
Why might closing prices be distorted on index rebalance days, and who benefits?
\end{interviewq}

\begin{interviewq}[$\star\star\star$ \iqroles{researcher}]\label{iq:hf:auctions-and-the-close:6}
With a linear impact curve $\Delta(I)=\lambda I$, a share $\pi$ that stays, and a provider offering $q$ at the reference, find the optimal $q$.
\end{interviewq}
